\section{The plan to Thursday}
\label{sec:plan}

Two days, three sittings and a mock. Each sitting has six core problems from White, in order, each a little harder than the last. Every one has its answer in the back of the book, so every call can be settled the same evening. Write the call line before each problem, as in \secref{sec:call}{reading a question}, and settle it as soon as you check the answer. Read the matching sections of this document before a sitting, not during it. Optional rungs are for a sitting that runs fast; the per-section ladders list more.

\subsection*{Tuesday 2026-11-10 evening: Chapter 3}

Read first: \secref{sec:momentum}{linear momentum} and \secref{sec:energy}{the energy equation}; skim the rest of Chapter 3.

\begin{center}\small
\begin{tabular}{@{}c l >{\raggedright\arraybackslash}p{5.6cm} >{\raggedright\arraybackslash}p{4.4cm}@{}}
\toprule
& Problem & What it is & What it trains \\
\midrule
1 & \white{P3.14} & A tank with one inlet and two outlets; the level rate, then the outlet speed for a constant level. & Unsteady mass, signs of inflow and outflow. \\
2 & \white{P3.40} & A round jet hits a fixed plate square on; the holding force. & The simplest momentum control volume. \\
3 & \white{P3.62} & A two-branch nozzle with angled jets; the bolt force at the inlet flange. & Pressure force at a cut, vector momentum flux. \\
4 & \white{P3.88} & A pump-jet boat with drag proportional to speed squared; the cruising speed. & Moving control volume, relative velocity. \\
5 & \white{P3.180} & A pump lifts water between reservoirs with losses and an efficiency; the power. & Energy equation in head form. \\
6 & \white{P3.150} & The moment about the inlet flange of an angled elbow. & Angular momentum, $\vect r \times \V$. \\
\midrule
opt & \white{P3.46} & A sheet splits on a frictionless inclined plate; the split fraction. & Zero tangential force as an assumption. \\
opt & \white{P3.122} & A tank filled steadily while draining through a hole; the limiting depth. & Bernoulli with unsteady mass. \\
\bottomrule
\end{tabular}
\end{center}

\subsection*{Wednesday 2026-11-11 morning: Chapter 4}

Read first: \secref{sec:navierstokes}{Navier--Stokes} and \secref{sec:exactflows}{exact viscous flows}. Before problem 1, write continuity and the Cartesian Navier--Stokes equations from memory, then check against the box in that section. Do it again before problem 4, in cylindrical coordinates.

\begin{center}\small
\begin{tabular}{@{}c l >{\raggedright\arraybackslash}p{5.6cm} >{\raggedright\arraybackslash}p{4.4cm}@{}}
\toprule
& Problem & What it is & What it trains \\
\midrule
1 & \white{P4.4} & A 2-D converging nozzle field; the accelerations and where the largest is. & The material derivative. \\
2 & \white{P4.22} & Axisymmetric flow with axial velocity linear in $z$; the general radial velocity. & Cylindrical continuity; the ``constant'' is a function. \\
3 & \white{P4.24} & A channel with porous walls; a velocity from continuity, the constant from the wall suction. & Continuity plus a boundary condition. \\
4 & \white{P4.80} & An oil film draining down a vertical wall; the profile. & Reducing Navier--Stokes; free-surface condition. \\
5 & \white{P4.88} & A rod sliding along the axis inside a fixed tube; the annular profile. & Cylindrical Navier--Stokes; a $\ln r$ solution. \\
6 & \white{P4.60} & A drain vortex with swirl falling as $1/r$; is it rotational, and the surface depth. & Vorticity, then Bernoulli where it is allowed. \\
\midrule
opt & \white{P4.92} & A tank emptying through a laminar pipe; the level history. & Poiseuille inside an unsteady mass balance. \\
opt & \white{P4.30} & The nozzle field of problem 1: does it satisfy Navier--Stokes, and the pressure field. & Pressure from a given velocity field. \\
\bottomrule
\end{tabular}
\end{center}

\subsection*{Wednesday 2026-11-11 afternoon: Chapter 5}

Read first: \secref{sec:pitheorem}{the Pi theorem} and \secref{sec:similarity}{modelling and similarity}.

\begin{center}\small
\begin{tabular}{@{}c l >{\raggedright\arraybackslash}p{5.6cm} >{\raggedright\arraybackslash}p{4.4cm}@{}}
\toprule
& Problem & What it is & What it trains \\
\midrule
1 & \white{P5.12} & Five variables give one group; show it is a ratio of familiar groups. & Dimensions from definitions. \\
2 & \white{P5.18} & Laminar flow in a triangular pore; the Pi form and the effect of doubling its size. & One group means a scaling law. \\
3 & \white{P5.38} & Spray droplet size against nozzle size, speed and liquid properties. & Repeating variables; Re and We appear. \\
4 & \white{P5.24} & Lift on a missile; the dimensional matrix, its rank, the groups. & Rank, and Mach as a group. \\
5 & \white{P5.80} & A ship and its towing-tank model under Froude scaling; speed, drag and power ratios. & Froude scaling laws. \\
6 & \white{P5.66} & A torpedo's cavitation speed at a new depth. & A similarity group with a pressure difference in it. \\
\midrule
opt & \white{P5.4} & A sphere tested in water scaled to a balloon in air. & Reynolds matching across fluids. \\
opt & \white{P5.62} & A small wind-turbine test scaled to a large machine. & Scaling speed and power together. \\
\bottomrule
\end{tabular}
\end{center}

\subsection*{Wednesday 2026-11-11 evening: the mock}

Ninety minutes, closed book, with only the crib sheet, at the exam's time if the day allows (18:10). Three call lines in the first five minutes, then the protocol of \secref{sec:examevening}{exam evening}.

\begin{enumerate}[nosep,leftmargin=*]
\item \white{P3.112}. An alcohol jet from a lossless nozzle hits a plate, and the holding force is known; the mass flow and the inlet pressure. Momentum then Bernoulli.
\item \white{P4.84}. A liquid film draining down the outside of a vertical rod; the equation, its boundary conditions, the profile and the flow rate. Cylindrical Navier--Stokes with a free surface. If it is too hard at the twenty-minute mark, switch to \white{P4.36}.
\item \white{P5.52}. A Pi analysis of a transition time in an accelerating pipe flow, then a scaling to a new case.
\end{enumerate}

Mark it against the book's answers that evening, and settle the three calls. Then read the Misses section of \texttt{readout.md}.

\subsection*{Thursday 2026-11-12}

No new problems. In the morning, the call drill from \secref{sec:call}{reading a question}: ten unsolved problems from the ladders, two minutes each, a call line for each. In the afternoon, the \emph{First moves} and \emph{Trap} boxes of every section, once. Then \secref{sec:examevening}{exam evening}.

\subsection*{A warning about the answer key}

A few answers in the back of this copy of White do not belong to the problems with the same numbers, apparently from a different printing: P3.6, P4.6, P5.14 and P5.22 among those checked. None of them is on a ladder. If an answer looks unrelated to its problem, it probably is; settle that call \emph{moot}.
